\clearpage
\VerdierPhysicalPageAnchor{0210}%
\begingroup\small\setlength{\parskip}{0pt}\setlength{\emergencystretch}{1.5em}
\setlength{\abovedisplayskip}{4pt plus 2pt minus 2pt}%
\setlength{\belowdisplayskip}{4pt plus 2pt minus 2pt}%

\VerdierSemanticAnchor{ch3d-3-1-4-proof}
証明すべきものは第三の主張だけである。他の主張は命題 (2.4.4) から直ちに
従う帰結である。主張 2) から、$\mathcal A$ のコホモロジー次元が、
$\mathcal A$ の対象の入射次元（それぞれ射影次元）の
上限以下であることが直ちに分かる。
$\mathcal A$ のコホモロジー次元が $n$ に等しいとし、

\VerdierSemanticAnchor{ch3d-3-1-4-proof-resolution}
\[
I^\bullet=\cdots\longrightarrow0\longrightarrow0\longrightarrow I^0
\longrightarrow I^1\longrightarrow I^2\longrightarrow I^3
\longrightarrow\cdots
\]
を $\mathcal A$ の対象 $Y$ の入射型分解とする。$\mathcal A$ の任意の対象
$X$ に対し、
$\operatorname{Hom}_{K(\mathcal A)}(X[-n-1],I^\bullet)=0$ である。
とくに $X=\operatorname{Ker}d_{I^\bullet}^{n+1}$ とすれば、標準射
$\operatorname{Ker}d_{I^\bullet}^{n+1}\longrightarrow I^{n+1}$ が
$I^n$ を経由して分解することが直ちに分かる：

\VerdierSemanticAnchor{ch3d-3-1-4-proof-factorization}
\[
\begin{tikzcd}[column sep=4.8em,row sep=2.8em]
{} & \operatorname{Ker}d_{I^\bullet}^{n+1}
  \arrow[dl,dashed] \arrow[d] & {} \\
I^n \arrow[r] & I^{n+1} \arrow[r] & \cdots
\end{tikzcd}
\]

\VerdierSemanticAnchor{ch3d-3-1-4-proof-truncation}
したがって $\operatorname{Ker}d_{I^\bullet}^{n+1}$ と
$\operatorname{Ker}d_{I^\bullet}^{n}$ は $I^n$ の直和因子であり、それゆえ
入射対象である。複体
\[
\cdots\longrightarrow0\longrightarrow I^0\longrightarrow I^1
\longrightarrow\cdots\longrightarrow I^{n-1}
\longrightarrow\operatorname{Ker}d_{I^\bullet}^{n}
\longrightarrow0\longrightarrow\cdots
\]
は $Y$ の入射分解である。したがって $Y$ の入射次元は $n$ 以下である。
同様に、$\mathcal A$ の任意の対象 $X$ の射影次元が $n$ 以下であることが
示される。

\VerdierSemanticAnchor{ch3d-3-1-5}
\medskip\noindent\textbf{3.1.5.}\quad
$X$ と $Y$ を $\mathcal A$ の二つの対象とする。命題 (3.1.4) は、圏
$\mathcal A$ が十分な入射対象または射影対象をもつとき、ここで定義した群
$\operatorname{Ext}_{\mathcal A}^{i}(X,Y)$ が古典的に定義される群と
一致することを示している。

\VerdierSemanticAnchor{ch3d-3-2}
\subsection*{3.2. Yoneda 構成。}

\VerdierSemanticAnchor{ch3d-3-2-1}
\noindent\textbf{3.2.1.}\quad
$X$ と $Y$ を $\mathcal A$ の二つの対象とする。
$E_{\mathcal A}^{n}(X,Y)$（$n\geq1$）を、完全列

\VerdierSemanticAnchor{ch3d-3-2-1-sequence}
\[
S^n(X,Y)=0\longrightarrow Y\longrightarrow Z_{n-1}
\longrightarrow Z_{n-2}\longrightarrow\cdots\longrightarrow Z_0
\xrightarrow{u_0}X\longrightarrow0.
\]
全体の集合とする。整数 $n$ を完全列 $S^n(X,Y)$ の長さという。
$E_{\mathcal A}^{n}(X,Y)$ の各元 $S^n(X,Y)$ に三つの複体を対応させる：
\endgroup
